\documentclass[10pt,a4paper]{amsart}
\usepackage{amsmath,amssymb,amsthm,graphicx}
% Retained minimal preamble for the Stacks Project chapter more-algebra.tex (More on Algebra).
% The chapter file pinned-original-source.tex carries only \input{preamble} on its first line, so no
% preamble text lives inside the chapter itself.  The Stacks Project uses one single shared preamble for
% every chapter; the snapshot stages it identically as sources/stacks-chow/preamble.tex and
% sources/stacks-etale-cohomology/preamble.tex.  The block below is copied mechanically, line for line,
% from sources/stacks-chow/preamble.tex lines 14-18 (the Xy-pic setup) and 154-217 (the native theorem
% declarations together with the mathematical macros).  Nothing was added, removed or redefined; the
% wrapper class and the amsmath/amssymb/amsthm/graphicx base packages are supplied by the assembler.
% Chapter: more-algebra.tex -- More on Algebra.
% For commutative diagrams we use Xy-pic
\usepackage[all]{xy}

% We use 2cell for 2-commutative diagrams.
\xyoption{2cell}

% Theorem environments.
%
\theoremstyle{plain}
\newtheorem{theorem}[subsection]{Theorem}
\newtheorem{proposition}[subsection]{Proposition}
\newtheorem{lemma}[subsection]{Lemma}

\theoremstyle{definition}
\newtheorem{definition}[subsection]{Definition}
\newtheorem{example}[subsection]{Example}
\newtheorem{exercise}[subsection]{Exercise}
\newtheorem{situation}[subsection]{Situation}

\theoremstyle{remark}
\newtheorem{remark}[subsection]{Remark}
\newtheorem{remarks}[subsection]{Remarks}

\numberwithin{equation}{subsection}

% Macros
%
\def\lim{\mathop{\mathrm{lim}}\nolimits}
\def\colim{\mathop{\mathrm{colim}}\nolimits}
\def\Spec{\mathop{\mathrm{Spec}}}
\def\Hom{\mathop{\mathrm{Hom}}\nolimits}
\def\Ext{\mathop{\mathrm{Ext}}\nolimits}
\def\SheafHom{\mathop{\mathcal{H}\!\mathit{om}}\nolimits}
\def\SheafExt{\mathop{\mathcal{E}\!\mathit{xt}}\nolimits}
\def\Sch{\mathit{Sch}}
\def\Mor{\mathop{\mathrm{Mor}}\nolimits}
\def\Ob{\mathop{\mathrm{Ob}}\nolimits}
\def\Sh{\mathop{\mathit{Sh}}\nolimits}
\def\NL{\mathop{N\!L}\nolimits}
\def\CH{\mathop{\mathrm{CH}}\nolimits}
\def\proetale{{pro\text{-}\acute{e}tale}}
\def\etale{{\acute{e}tale}}
\def\QCoh{\mathit{QCoh}}
\def\Ker{\mathop{\mathrm{Ker}}}
\def\Im{\mathop{\mathrm{Im}}}
\def\Coker{\mathop{\mathrm{Coker}}}
\def\Coim{\mathop{\mathrm{Coim}}}

% Boxtimes
%
\DeclareMathSymbol{\boxtimes}{\mathbin}{AMSa}{"02}

%
% Macros for moduli stacks/spaces
%
\def\QCohstack{\mathcal{QC}\!\mathit{oh}}
\def\Cohstack{\mathcal{C}\!\mathit{oh}}
\def\Spacesstack{\mathcal{S}\!\mathit{paces}}
\def\Quotfunctor{\mathrm{Quot}}
\def\Hilbfunctor{\mathrm{Hilb}}
\def\Curvesstack{\mathcal{C}\!\mathit{urves}}
\def\Polarizedstack{\mathcal{P}\!\mathit{olarized}}
\def\Complexesstack{\mathcal{C}\!\mathit{omplexes}}
% \Pic is the operator that assigns to X its picard group, usage \Pic(X)
% \Picardstack_{X/B} denotes the Picard stack of X over B
% \Picardfunctor_{X/B} denotes the Picard functor of X over B
\def\Pic{\mathop{\mathrm{Pic}}\nolimits}
\def\Picardstack{\mathcal{P}\!\mathit{ic}}
\def\Picardfunctor{\mathrm{Pic}}
\def\Deformationcategory{\mathcal{D}\!\mathit{ef}}
\usepackage{cleveref}
\crefname{originallemma}{original lemma}{original lemmas}
\Crefname{originallemma}{Original lemma}{Original lemmas}
\makeatletter
\newlabel{more-algebra-lemma-regular-koszul-regular}{{062F}{source}{}{sourcepointer-1}{}}
\newlabel{more-algebra-lemma-regular-koszul-regular@cref}{{[originallemma][0][]062F}{[1][0][]source}{}{}{}}
\makeatother
\crefname{originallemma}{original lemma}{original lemmas}
\Crefname{originallemma}{Original lemma}{Original lemmas}
\makeatletter
\newlabel{lemma-cone-koszul}{{0629}{source}{}{sourcepointer-2}{}}
\newlabel{lemma-cone-koszul@cref}{{[originallemma][0][]0629}{[1][0][]source}{}{}{}}
\makeatother
\title{An M-regular sequence is M-Koszul-regular}
\author{The Stacks Project Authors}
\date{stacks-project:more-algebra.tex:local-snapshot-20261012; GNU FDL 1.2 or later}
\begin{document}
\maketitle
\setcounter{subsection}{1}
\noindent\textit{Editorial note.} All delivered mathematical text is the author's own, copied from the pinned Stacks Project chapter file more-algebra.tex; no statement or proof step was rewritten, reordered or removed.\par\medskip
\noindent\textit{Editorial note.} The chapter file pinned-original-source.tex carries only \textbackslash{}input\{preamble\} on its first line, so no preamble text lives inside the chapter. The retained file sources/stacks-more-algebra/retained-preamble.tex is a mechanical, line-for-line copy of the project's single shared preamble as staged in the snapshot (Xy-pic setup plus theorem declarations and mathematical macros). Nothing was added, removed or redefined, and the excerpt serves only as preamble, never as a span coordinate.\par\medskip
\noindent\textit{Editorial note.} Both counted spans start at the \textbackslash{}begin line of their environment, so the statement and the author's proof are delivered whole; surrounding prose of the section is deliberately outside the counted spans and is carried as standard prerequisites instead.\par\medskip
\noindent\textit{Editorial note.} The citation fields (source\_id, source\_version, authors, source\_url, license.name, source\_title) are identical to those of every other card this lane has produced from this chapter, so the batch selector treats them as one source.\par\medskip
\noindent\textit{Allowed standard prerequisites (editorial).}\begin{itemize}
\item Koszul complex of a sequence relative to a module and its homology, and the definition of M-Koszul-regularity. used throughout the proof, lines 7273-7298. Supplies the goal and the complex through which it is tested.
\item The Koszul complex of a longer sequence is the cone on multiplication by the first element on the Koszul complex of the tail. the author cites Lemma lemma-cone-koszul (tag 0629) at proof lines 7281-7284. Is the structural statement that converts the problem into a statement about the cone.
\item The long exact sequence of cohomology attached to a cone and the acyclicity criterion for a cone. used at proof lines 7284-7297. Transfers the vanishing of the homology from the tail to the whole sequence.
\end{itemize}\medskip
\ifdefined\hypertarget\hypertarget{sourcepointer-1}{}\fi
\noindent\textit{Source reference.} Original lemma 062F: focused-seeds/AGP/sources/stacks-snapshot-20261011/tags.txt, source lines 7794--7794. The lemma's own canonical tag, quoted so that readers can resolve the label inside the delivered text; it points at the counted statement itself and carries no proof content.\par\medskip
\ifdefined\hypertarget\hypertarget{sourcepointer-2}{}\fi
\noindent\textit{Source reference.} Original lemma 0629: focused-seeds/AGP/sources/stacks-snapshot-20261011/tags.txt, source lines 7788--7788. Navigation pointer to a statement the counted proof refers to but does not contain; nothing load-bearing is delegated to it.\par\medskip
\begin{lemma}
\label{lemma-regular-koszul-regular}
Let $R$ be a ring, $M$ an $R$-module, and $f_1, \ldots, f_r \in R$
such that for $i = 1, \ldots, r$ multiplication by $f_i$
is injective on $M/(f_1, \ldots, f_{i - 1})M$. Then $f_1, \ldots, f_r$
is $M$-Koszul regular. In particular, an $M$-regular sequence
is $M$-Koszul-regular and any regular sequence is Koszul-regular.
\end{lemma}

\begin{proof}
Let $R$, $M$, $f_1, \ldots, f_r$ be as in the first sentence of the lemma.
If $r = 1$, it is immediate that $f_1$ is $M$-Koszul-regular.
Assume $r > 1$. Since $f_1$ is a nonzerodivisor on $M$, we obtain
a short exact sequence of complexes:
$$
0 \to K_\bullet(f_2, \ldots, f_r) \otimes M
\xrightarrow{f_1}
K_\bullet(f_2, \ldots, f_r) \otimes M \to
K_\bullet(\overline{f}_2, \ldots, \overline{f}_r) \otimes M/f_1M \to 0
$$
Here $\overline{f}_i$ is the image of $f_i$ in $R/(f_1)$. By
Lemma \ref{lemma-cone-koszul}
the complex $K_\bullet(f_1, \ldots, f_r)$
is isomorphic to the cone of multiplication by $f_1$
on $K_\bullet(f_2, \ldots, f_r)$. Thus
$K_\bullet(R, f_1, \ldots, f_r) \otimes M$ is isomorphic
to the cone on the first map. Hence
$K_\bullet(\overline{f}_2, \ldots, \overline{f}_r) \otimes M/f_1M$
is quasi-isomorphic to $K_\bullet(f_1, \ldots, f_r) \otimes M$.
As $R/(f_1)$, $M/f_1M$, $\overline{f}_2, \ldots, \overline{f}_r$
satisfy the conditions of the lemma, by induction we conclude
this complex is acyclic in postive degrees. This finishes the proof
of the first statement. The second statement immediately follows
from the first.
\end{proof}

\end{document}
